\section{Methodology: From Algebraic to Geometric Tracking}

To address the limitations of purely algebraic tracking while maintaining real-time performance on edge devices, we design a three-stage tracking architecture that progressively incorporates hardware-aware acceleration and geometric reasoning. Each stage builds upon the previous one, transitioning from efficient 2D correspondence estimation to robust 6-DoF pose refinement on $\mathrm{SE}(3)$.

\begin{figure}[p]
\centering
\resizebox{1\textwidth}{!}{%
\begin{tikzpicture}[x=1cm, y=1cm]

%% ═══════════════════════════════════════════════════════════════════════
%%  QUERY STACK  (top, centred)
%% ═══════════════════════════════════════════════════════════════════════
\foreach \i in {\numscales,...,1} {
    \pgfmathsetmacro{\dx}{(\i-1)*\deltadim}
    \pgfmathsetmacro{\dy}{(\i-1)*\deltadim}
    \node[querynode, anchor=north west,]
        (render\i) at ($(\dx,\dy)$) {Render Mesh};
    \node[querynode, anchor=north west,]
        (coords\i) at (render\i.north east) {Coord Buffer};
    \node[querynode, anchor=north west,]
        (feats\i)  at (coords\i.north east) {GPU Descriptors};
}
\node[font=\scriptsize\sffamily\itshape, text=\querydraw,
      right=0.28cm of feats1, align=left] {$L$\\levels};

\draw[fw] ($(render1.west)+(-3.0,0)$) -- (render1.west)
    node[ann, pos=0.5, above] {Query Image/Mesh};

%% ═══════════════════════════════════════════════════════════════════════
%%  CAMERA + DOWNSCALE + STATE MACHINE
%% ═══════════════════════════════════════════════════════════════════════
\node[camnode, anchor=north, below=2.0cm of coords1] (DS)
    {Downscale \& Gray\\[-1pt]\scriptsize $(\leq R_{\max}$ px$)$};
\draw[fw] ($(DS.west)+(-2.8,0)$) -- (DS.west)
    node[ann, pos=0.4, above] {Camera Frame};

\node[smnode, below=0.65cm of DS] (SM)
    {\textbf{State Machine}\\[-2pt]
     \scriptsize \textsc{matching} $\;\big|\;$ \textsc{optical flow}};
\draw[fw] (DS.south) -- (SM.north);

\coordinate (Lcol) at ($(SM.south)+(-3.4,0)$);
\coordinate (Rcol) at ($(SM.south)+(+3.4,0)$);
\coordinate (Mcol) at ($(SM.south)+(0,0)$);

%% ═══════════════════════════════════════════════════════════════════════
%%  LEFT BRANCH  –  MATCHING
%% ═══════════════════════════════════════════════════════════════════════
\node[matchnode] (GF) at ($(Lcol)+(0,-1.55)$)
    {GPU Feature Extraction\\[-1pt]
     \scriptsize Shi-Tomasi $\!\to\!$ orientation $\!\to\!$ binary desc};

\draw[fw, rounded corners=4pt]
    (SM.south) -- ++(0,-0.25)
    node[ann, left=3pt] {\textsc{matching}}
    -- ++(0,-0.35) -| (GF.north);

\node[matchnode, below=0.48cm of GF] (MS)
    {Multiscale Match\\[-1pt]
     \scriptsize $\ell^*\!=\!\arg\max_\ell|\mathcal{M}_\ell|$\enspace Lowe ratio $\tau$};
\draw[fw] (GF.south) -- (MS.north);

\draw[qfeed, rounded corners=3pt]
    (feats1.south)
    -- ++(0,-0.6)
    -- ++(-10.2,0)
    |- (MS.west);

\node[matchnode, below=0.48cm of MS] (HG)
    {Homography (RANSAC)\\[-1pt]
     \scriptsize reprojection $<3$px · quad check};
\draw[fw] (MS.south) -- (HG.north);

%% ═══════════════════════════════════════════════════════════════════════
%%  RIGHT BRANCH  –  OPTICAL FLOW
%% ═══════════════════════════════════════════════════════════════════════
\node[lknode] (LK) at ($(Rcol)+(0,-1.55)$)
    {Lucas--Kanade Tracking\\[-1pt]
     \scriptsize prev $\!\to\!$ curr frame · pyramidal LK};

\draw[lkarr, rounded corners=4pt]
    (SM.south) -- ++(0,-0.25)
    node[ann, right=3pt] {\textsc{optical flow}}
    -- ++(0,-0.35) -| (LK.north);

\node[lknode, below=0.48cm of LK] (FBV)
    {Fwd--Bwd Validation\\[-1pt]
     \scriptsize $\|\mathbf{x}_{t-1}\!-\!\hat{\mathbf{x}}_{t-1}\|<\varepsilon$
     · boundary check};
\draw[lkarr] (LK.south) -- (FBV.north);

\node[lknode, below=0.48cm of FBV] (LKHG)
    {Homography (RANSAC)\\[-1pt]
     \scriptsize LK pts vs.\ query pts · quad check};
\draw[lkarr] (FBV.south) -- (LKHG.north)
    node[ann, pos=0.5, right=3pt] {status pts only};

%% ═══════════════════════════════════════════════════════════════════════
%%  MERGE — 2D–3D Lifting + SE(3) + Output
%% ═══════════════════════════════════════════════════════════════════════
\coordinate (mergeY) at ($(Mcol |- LKHG.south)+(0,-0.55)$);

\node[posenode, anchor=north] (LIFT) at (mergeY)
    {2D--3D Lifting\\[-1pt]
     \scriptsize inliers only · coord buffer at $\ell^*$};

\draw[fw, rounded corners=4pt]
    (HG.south) -- ++(0,-0.2)
    -| node[ann, pos=0.6, right=2pt] {inliers}
    (LIFT.north west);

\draw[lkarr, rounded corners=4pt]
    (LKHG.south) -- ++(0,-0.2)
    -| node[ann, pos=0.6, left=2pt] {inliers}
    (LIFT.north east);

\draw[qfeed, rounded corners=3pt]
    (coords1.south)
    -- ++(0,-1)
    -- ++(-6.2,0)
    |- (LIFT.west);

\node[posenode, below=0.5cm of LIFT] (SE3)
    {$\mathrm{SE}(3)$ Gauss--Newton Refinement\\[-2pt]
     \scriptsize $\mathbf{T}\!\leftarrow\!\exp(\hat{\boldsymbol{\xi}})\mathbf{T}$
     \enspace
     $\min_{\boldsymbol{\xi}}\!\sum_i\|\mathbf{x}_i\!-\!\pi(\mathbf{KTX}_i)\|^2$};
\draw[pose] (LIFT.south) -- (SE3.north);

\node[smnode, below=0.5cm of SE3] (TR)
    {Pose Output · confirm\texttt{++}\\[-1pt]
     \scriptsize $\geq\!$\texttt{MIN\_CONFIRM\_FRAMES} $\Rightarrow$
     tracking $=$ \texttt{true} · \textsc{optical flow}};
\draw[pose] (SE3.south) -- (TR.north);

%% ═══════════════════════════════════════════════════════════════════════
%%  FAIL NODES
%% ═══════════════════════════════════════════════════════════════════════
\node[warnnode, anchor=center]
    (FAIL) at ($(HG.west)+(-2.55, -1)$)
    {Re-Match\strut\\[-2pt]quad fail $|$ inliers${<}\theta$ $|$ matches${<}\theta$};

\draw[warn] (HG.west) --
    node[annW, pos=0.5, above=1pt, rotate=90] {fail} (FAIL);
\draw[warn] (MS.west) --
    node[annW, pos=0.5, above=1pt, rotate=90] {matches${<}\theta$} (FAIL);
\draw[warn, rounded corners=4pt]
    (FAIL.north) -- ++(-0.3, 0)
    |- node[annW, pos=0.2, above=1pt, rotate=90] {tracking $=$ \texttt{false}}
    (GF.west);

\node[warnnode, anchor=center, rotate=180]
    (OFWEAK) at ($(LKHG.east)+(1.55, 0.5)$)
    {OF Weak / Lost\strut\\[-2pt]pts${<}\texttt{LK\_MIN}$ $|$ quad fail};

\draw[warn, rounded corners=3pt]
    (LKHG.east) -- node[annW, pos=0.5, above=1pt, rotate=-90] {fail} (OFWEAK);
\draw[warn, rounded corners=3pt]
    (FBV.east) --
    node[annW, pos=0.3, above=1pt, rotate=-90] {pts${<}\theta$} (OFWEAK);
\draw[warn, rounded corners=4pt]
    (OFWEAK.north) -- ++(0.3,0)
    |- node[annW, pos=0.2, above=1pt, rotate=-90] {tracking $=$ \texttt{false} · reset}
    (LK.east);

\begin{pgfonlayer}{background}
\node[
    fill=gray!10,
    draw=gray!60!black,
    rounded corners=8pt,
    inner sep=24pt,
    fit=(DS)(SM)(GF)(MS)(HG)(LK)(FBV)(LKHG)(LIFT)(SE3)(TR)(OFWEAK)(FAIL)
] {};
\end{pgfonlayer}

\end{tikzpicture}
}
\caption{%
Overview of the GPU-centric detect--refine--track framework.
\textbf{Top:} A stack of $L$ rendered reference meshes, object-space coordinates, and GPU descriptors.
\textbf{Centre:} The camera frame is downscaled and grayscaled, then fed into a state machine that selects between global matching (left) or optical-flow-based tracking (right).
\textbf{Left branch (matching):} GPU feature extraction, multiscale matching, and homography verification (RANSAC) produce inliers for 2D--3D lifting.
\textbf{Right branch (optical flow):} Lucas--Kanade tracking with forward--backward validation and homography verification produces inliers for 2D--3D lifting.
\textbf{Merge:} Inliers from both branches feed 2D--3D lifting, followed by $\mathrm{SE}(3)$ Gauss--Newton refinement and pose output confirmation.
\textbf{Fail handling:} Re-match or optical-flow failure nodes loop back to the corresponding branches when matches or confidence fall below thresholds.
}
\label{fig:pipeline}
\end{figure}


\subsection{Stage 1: Optimized CPU-Based Tracking (Algebraic)}

The first stage consists of a feature-based tracking pipeline executed entirely on the CPU. This stage relies on 2D algebraic constraints for motion estimation and is carefully optimized for high throughput on mobile processors.

\paragraph{Feature Extraction and Matching.}
Incoming frames are converted to grayscale, and keypoints are detected using a high-speed corner detector (e.g., FAST), followed by non-maximum suppression to enforce spatial sparsity. For each keypoint, orientation is estimated using local image gradients, and binary descriptors (e.g., BRISK) are constructed. Descriptor matching is performed against a precomputed reference set using Hamming distance.

\paragraph{Algebraic Pose Initialization.}
Given a set of 2D correspondences $\{\mathbf{x}_i \leftrightarrow \mathbf{x}'_i\}$, a homography $\mathbf{H}$ is estimated using Direct Linear Transform (DLT) with RANSAC for outlier rejection:
\begin{equation}
\mathbf{x}'_i \sim \mathbf{H}\mathbf{x}_i.
\end{equation}
While efficient, this formulation assumes planar structure and becomes unstable under strong perspective distortion, scale changes, or out-of-plane rotations.

\paragraph{Optical Flow Tracking.}
After initialization, tracking proceeds using a pyramidal Lucas--Kanade (LK) optical flow method to propagate inlier correspondences across frames, reducing the need for repeated feature extraction.

\paragraph{SIMD Optimization.}
All components in this stage are optimized using SIMD vectorization (e.g., ARM NEON) and compiler auto-vectorization to maximize data-level parallelism. Despite these optimizations, feature extraction and matching remain a dominant computational bottleneck.

\subsection{Stage 2: GPU-Accelerated Feature Processing (Algebraic)}

To alleviate CPU bottlenecks, the second stage offloads feature extraction and descriptor computation to the GPU.

Grayscale conversion, keypoint detection, and descriptor construction are implemented using custom fragment shaders, enabling per-pixel parallel processing across the image domain. Descriptor bits are generated via pairwise intensity comparisons within rotated sampling patterns directly in the shader pipeline. Descriptor matching is further parallelized using compute shaders, significantly reducing matching latency.

The resulting correspondences are transferred back to the CPU for homography estimation and optical flow tracking. While this stage substantially improves throughput, the underlying motion model remains algebraic and inherits the same geometric limitations as Stage 1.

\subsection{Stage 3: Geometry-Aware $\mathrm{SE}(3)$ Render-and-Track Refinement}

To overcome the instability of purely algebraic methods, the third stage introduces a geometry-aware pose estimation framework that operates directly on the $\mathrm{SE}(3)$ manifold.

\paragraph{Depth-Augmented Reference Modeling.}
To resolve scale ambiguity, the known object geometry is rendered at a discrete set of depth hypotheses. For each hypothesis, two buffers are generated:
\begin{itemize}
    \item \textbf{Appearance buffer:} Used for feature extraction and matching.
    \item \textbf{Coordinate buffer:} Stores per-pixel object-space coordinates of visible surface points.
\end{itemize}

\paragraph{Direct 2D--3D Correspondence Lifting.}
Each valid descriptor match directly retrieves its corresponding 3D point $\mathbf{X}_i$ from the coordinate buffer, forming a set of 2D--3D correspondences without requiring explicit back-projection or triangulation.

\paragraph{Tension-Based Pose Refinement on $\mathrm{SE}(3)$.}

Let $\mathbf{x}_i$ denote the observed normalized image coordinate and $\mathbf{p}_i(\mathbf{T}) = \pi(\mathbf{K}\mathbf{T}\mathbf{X}_i)$ denote the projected point under pose $\mathbf{T} \in \mathrm{SE}(3)$. The reprojection residual is:
\begin{equation}
\mathbf{e}_i = \mathbf{x}_i - \mathbf{p}_i(\mathbf{T}), \quad d_i = \|\mathbf{e}_i\|.
\end{equation}

We model each correspondence as a virtual spring in image space, where the residual magnitude $d_i$ defines the spring length. The smallest residual
\[
L = \min_i d_i
\]
is treated as a reference equilibrium. The resulting tension induces an adaptive weight:
\begin{equation}
w_i = \frac{d_i - L}{d_i}.
\end{equation}

This formulation assigns higher influence to inconsistent correspondences while suppressing those already close to agreement.

Pose refinement is then formulated as a weighted nonlinear least-squares problem:
\begin{equation}
E(\mathbf{T}) =
\sum_i w_i \left\| \mathbf{x}_i - \pi(\mathbf{K}\mathbf{T}\mathbf{X}_i) \right\|^2.
\end{equation}

Linearizing on the Lie algebra $\mathfrak{se}(3)$ yields the weighted Gauss--Newton update:
\begin{equation}
(\mathbf{J}^\top \mathbf{W} \mathbf{J})\,\boldsymbol{\delta}
=
\mathbf{J}^\top \mathbf{W} \mathbf{e},
\end{equation}
where $\mathbf{W} = \mathrm{diag}(w_i)$.

The pose is updated via the exponential map:
\begin{equation}
\mathbf{T} \leftarrow \exp(\boldsymbol{\delta}^\wedge)\mathbf{T}.
\end{equation}

This physics-inspired formulation adaptively reweights correspondences based on relative residual structure, improving robustness without requiring manually tuned robust loss functions.

\paragraph{System Integration.}
The complete system operates as a detect--refine--track pipeline. Global feature matching provides initialization, $\mathrm{SE}(3)$ refinement ensures geometric consistency, and optical flow enables efficient frame-to-frame tracking. This hybrid design achieves robust 6-DoF pose estimation while maintaining real-time performance on mobile-class hardware.